\section{What the Boundary Means}
\label{sec:analysis}

\paragraph{Capability is not complementarity.}
Terra answers more effectively under the iterative policy than the other two
tiers: its ReAct score exceeds control ReAct by 0.063
($[0.002,0.124]$), while its Static score does not differ reliably from control.
This produces the positive interaction and answers the principal robustness
objection: ReAct's earlier loss was partly model-dependent. But Terra's ReAct
successes almost entirely overlap its Static successes. The system-level tie
therefore does not imply a useful mixture.

The distinction is mechanical. A selector can recover at most the oracle union,
whose gain over the better branch equals the rarer one-sided cell. Here that
cell contains one or two items per tier. A broad router would send many cases to
a branch with little unique value; a narrow router would have too few benefit
examples to learn reliably. Calls, latency, and cost all increase for ReAct, so
there is no compensating efficiency frontier.

\paragraph{Why contract fairness matters.}
The pilot used unequal final-answer and context contracts, and a post hoc author
audit found scorer-sensitive one-sided outputs. We neither rescore those cases
nor call the audit independent adjudication. The harmonized and capability
studies instead share answer semantics, ceilings, scorer input, one-attempt
rules, and provider binding. They share a retrieval API and budget, not an
identical retrieved dossier: Static averages about 2.85 retrieval operations,
whereas ReAct adaptively uses 1.16--1.31. The estimates therefore compare
deployed retrieval policies, including query and evidence-selection effects.

The practical sequence is: harmonize evaluation contracts, measure paired
complementarity, and only then optimize a selector. Retuning on the present
outcomes and reporting the same sample as evaluation would invalidate that
sequence.
